\section{Conclusão}
\label{sec:conclusao}

Apresentamos um simulador de fluxo Kanban com cadência Scrum, sinergia par a par e traços individuais, com foco em \textbf{transparência algorítmica} e \textbf{reprodutibilidade} via semente de PRNG. O texto combina formalização (Seção~\ref{sec:modelo}), calendário de ritos (Seção~\ref{sec:ritos}), arquitetura do protótipo (Seção~\ref{sec:arquitetura}) e um pacote mínimo de cenários A/B com mesma carga inicial e mesmos parâmetros globais, permitindo discutir sensibilidade de métricas de fluxo a fatores interpessoais operacionalizados de forma explícita.

Em síntese, a contribuição central não é ``prever'' um time real, mas disponibilizar um \textbf{laboratório computacional} com fronteira clara entre hipótese operacional e evidência empírica. Como extensões, citamos: calibração com dados de projetos, catálogo de traços mais rico (com testes de face validity), políticas alternativas de priorização, e integração com protocolos de experimentação com equipes humanas jogando o simulador.
